\subsection{Data Resources}
\label{sec:data}

The system consumes two complementary data resources: a plant neutron-current measurement archive that supplies observable signal phenomena, and a technical evidence corpus that supplies retrievable documentary support. Their roles are deliberately distinct: neither resource alone is treated as a fault-label database.

\subsubsection{Neutron-current measurement archive}
The measurement archive contains 42 multi-channel neutron-current CSV records from four instrumentation-bank groups (A1, A2, B1, and B2) in the same plant campaign. Each record contains one time column and seven neutron-current channels (eight columns overall). The archive composition is summarized in \cref{tab:neutron-data}.

\begin{table}[H]
  \centering
  \caption{Composition of the neutron-current measurement archive.}
  \label{tab:neutron-data}
  \small
  \resizebox{\linewidth}{!}{%
  \begin{tabular}{@{}lccc@{}}
    \toprule
    Group & Number of series & Channels per series & Columns per series \\
    \midrule
    A1 & 11 & 7 & 8 \\
    A2 & 11 & 7 & 8 \\
    B1 & 10 & 7 & 8 \\
    B2 & 10 & 7 & 8 \\
    \midrule
    Total & 42 & --- & --- \\
    \bottomrule
  \end{tabular}
  }
\end{table}

Channel identifiers follow plant tagging conventions (e.g., \texttt{RIC}\ldots\texttt{MA\_VER}). Within a group, the seven detectors occupy a circumferential layout, so their ordering in the CSV is associated with relative channel position rather than an arbitrary list. Representative series contain on the order of $6\times10^{5}$ samples at a nominal \SI{1}{\second} interval over multi-day windows.

The raw exports can contain local timestamp disorders and non-monotonic segments. Chronological sorting is therefore a mandatory first operation. The present workflow does not use a uniform sample-level fault ground truth: its output is a condition report of observed signal events, not a supervised fault class. Derived reports are workflow artifacts rather than annotations attached to the raw CSV records.

\subsubsection{Technical evidence corpus}
The evidence resource is a curated 64-document collection of nuclear instrumentation and control manuals converted from PDF to Markdown and accompanied by document-source metadata. It contains a 13-document IAEA core, including guidance on online monitoring and instrument-channel assessment~\citep{iaea2008olm1,iaea2008olm2}, together with 51 same-domain supplementary documents from sources such as NRC, EPRI/DOE, OECD/NEA, and laboratory or topical reports. The historical and CPU protocols use distinct index and chunk constructions, as detailed in \cref{sec:eval}, but draw from this same document collection.

The supplementary documents are not assumed to be irrelevant. They enlarge the searchable knowledge base, while the IAEA core represents a deployment policy that prioritizes a pre-designated authority set. This distinction motivates the source-priority measurements in the evaluation but does not redefine source priority as universal relevance.

\subsubsection{Data roles and semantic boundary}
The two resources meet only through an auditable semantic interface. Deterministic probes first convert a measurement record $\mathbf{X}$ into observed events and a structured alert summary $\mathbf{a}=g(\mathbf{X})$. Retrieval then converts $\mathbf{a}$ into terminology-rich queries over the evidence corpus. Thus a measurement anomaly is not directly promoted to a physical diagnosis, and a manual passage is not treated as a direct explanation of an unprocessed waveform. The resulting evidence-linked memo remains an aid to engineering review.
